%% ----------------------------------------------------------------
%% Gaussian integrals and Gaussian processes problem sheet
%% (lectures: gaussianIntegrals, gaussianProcesses)
%% ----------------------------------------------------------------
\documentclass{sotonExamBoxes}    % use option [answers] to include them
%% ----------------------------------------------------------------
\usepackage{graphicx}
\usepackage{alltt}
\usepackage{bm}
\usepackage{amssymb}
\usepackage{amsfonts}
\usepackage{amsmath}
\thinmuskip=1mu
%% macros.tex -- shared macro definitions for the problem sheets.
%%
%% Every sheet inputs this instead of carrying its own copy. Input it after
%% the \usepackage lines: the definitions below use bm, amsmath and amssymb.
%%
%% Six macros were defined inconsistently across the sheets when this file
%% was made; the choices are marked CONFLICT below.




\usepackage{stmaryrd}

\renewcommand{\ttdefault}{pcr}

\newcommand{\tr}{\textsf{T}}

\newcommand{\e}[1]{{\rm e}^{#1}}

\newcommand{\E}[1]{{\rm e}^{{\displaystyle #1}}}

\newcommand{\logg}[1]{\log\left(#1\right)}

\newcommand{\Prob}[1]{\mathbb{P}\left(#1\right)}

\newcommand{\Step}{\mathop{\mathrm{O}\hspace{-15pt}\rule[8.5pt]{11pt}{1.2pt}\hspace{2pt}}}

\newcommand\diag{\mathop{\mathrm{diag}}}

\newcommand{\grad}{\bm{\nabla}}

\newcommand{\dd}{\mathrm{d}}

\newcommand{\class}{\mathcal{C}}

\newcommand{\data}{\mathcal{D}}

%% CONFLICT: gaussianProcess.tex had {U}{eur}{b}{n} (Euler bold);
%% every other sheet had this. Matrices in that one sheet now look
%% different -- see the report.

\DeclareMathAlphabet{\mat}{OT1}{cmss}{bx}{n}

%% The identity needs to be distinguishable. In cmss a bold I is a plain
%% vertical bar with no crosspieces, so \mat{I} reads as a rule rather than a
%% letter. Euler's bold I keeps its serifs, so the identity alone is set in
%% Euler while everything else stays cmss.
\DeclareMathAlphabet{\matEuler}{U}{eur}{b}{n}
\newcommand{\Id}{\matEuler{I}}

%% CONFLICT and BUG FIX: 16 sheets defined this using \qpartial,
%% which is defined nowhere, so \pd was broken in all of them.
%% Only problem2.tex had the working \partial form, kept here.

\newcommand{\pd}[2]{\ensuremath{\frac{\partial #1}{\partial #2}}\xspace}

\newcommand{\M}[1]{\ensuremath{\boldsymbol{#1}}\xspace}

\newcommand{\bst}{\rule{0mm}{5mm}}

%% CONFLICT: some sheets omitted the \! . Kept the tighter spacing,
%% which is what every sheet that actually uses \av had.

\newcommand{\av}[2][\,]{\mathbb{E}_{#1}\!\left[ {\strut #2} \right]}

\newcommand{\len}[1]{\| #1 \|}

\newcommand{\inner}[2]{\left\langle #1, #2\right\rangle}

\newcommand{\squeeze}{\setlength{\itemsep}{-2pt}}

%% CONFLICT: two spacings, identical output in maths mode.

\newcommand{\bra}[1]{\left( #1 \right)}

\newcommand{\pred}[1]{\left\llbracket { \small #1} \right\rrbracket}

\newcommand{\mln}{\hat{f}_n(\bm{x})}


\newcommand{\Tr}{\mathop{\mathrm{tr}}\,}

%% NOT SHARED: gaussianProcess.tex redefines \det to take an argument
%% (\det{\mat{V}}), while calculus.tex and vectorSpaces.tex use the standard
%% no-argument operator (\det(\mat{X})). The two cannot coexist, so that
%% redefinition stays local to gaussianProcess.tex.

%% NOT SHARED: \ml and \mm mean different things in different sheets.
%% amgm and convexity use a no-argument \ml and \mm = \hat{m}_c(x);
%% ensembleLearning uses a one-argument \ml and \mm = \hat{m}(x).
%% Both forms are in use, so each sheet keeps its own.

%% local: this sheet writes \det{X}, not the standard \det X
\renewcommand{\det}[1]{\textrm{det}({#1})}
\graphicspath{{figures/}}
\begin{document}
\unitTitle{Gaussian Integrals and Gaussian Processes Problem Sheet}
\authors{Adam Pr\"ugel-Bennett}
\unitcode{AICE3001/AICE6002}
\semester{2}
\year{2022}
\duration{45}{0}

\maketitle

This sheet goes with the lectures on Gaussian integrals and Gaussian
processes.  Questions~1 and~2 only ask you to \emph{use} the results
below; Questions~3 and~4 derive them, and are longer and harder.  Do
Questions~1 and~2 first.

\medskip
\noindent\fbox{\begin{minipage}{0.96\textwidth}
\textbf{The Gaussian toolkit}
\begin{itemize}
\item \textbf{Integrals:} $\displaystyle\int_{-\infty}^{\infty}\e{-\frac{a}{2}x^2+b\,x}\,\dd x
  = \sqrt{\frac{2\pi}{a}}\;\e{\frac{b^2}{2a}}$ for $a>0$, and in $n$
  dimensions, for positive definite $\mat{A}$,
  $\displaystyle\int\e{-\frac12\bm{x}^\tr\mat{A}\bm{x}+\bm{b}^\tr\bm{x}}\,\dd\bm{x}
  = \sqrt{\frac{(2\pi)^n}{\det{\mat{A}}}}\;\e{\frac12\bm{b}^\tr\mat{A}^{-1}\bm{b}}$.
\item \textbf{Moment generating function:} if $\bm{x}\sim\mathcal{N}(\bm{\mu},\mat{\Sigma})$
  then $\av{\e{\bm{t}^\tr\bm{x}}} = \e{\bm{t}^\tr\bm{\mu}+\frac12\bm{t}^\tr\mat{\Sigma}\bm{t}}$
  (in one dimension $\av{\e{tX}} = \e{\mu t+\frac12\sigma^2t^2}$).
\item \textbf{Sums:} if $X\sim\mathcal{N}(\mu_1,\sigma_1^2)$ and
  $Y\sim\mathcal{N}(\mu_2,\sigma_2^2)$ are independent then
  $X+Y\sim\mathcal{N}(\mu_1+\mu_2,\sigma_1^2+\sigma_2^2)$.
\item \textbf{Linear maps:} $\mat{W}\bm{x}+\bm{c}\sim
  \mathcal{N}(\mat{W}\bm{\mu}+\bm{c},\;\mat{W}\mat{\Sigma}\mat{W}^\tr)$.
\item \textbf{Marginals and conditionals:} for $\begin{pmatrix}\bm{z}_1\\ \bm{z}_2\end{pmatrix}\sim
  \mathcal{N}\!\left(\begin{pmatrix}\bm{\mu}_1\\ \bm{\mu}_2\end{pmatrix},
  \begin{pmatrix}\mat{\Sigma}_{11}&\mat{\Sigma}_{12}\\
    \mat{\Sigma}_{21}&\mat{\Sigma}_{22}\end{pmatrix}\right)$:
  \begin{itemize}
  \item marginal: $\bm{z}_1\sim\mathcal{N}(\bm{\mu}_1,\mat{\Sigma}_{11})$
    (strike out the rows and columns of $\bm{z}_2$);
  \item conditional: $\bm{z}_1\mid\bm{z}_2\sim\mathcal{N}\!\left(
    \bm{\mu}_1+\mat{\Sigma}_{12}\mat{\Sigma}_{22}^{-1}(\bm{z}_2-\bm{\mu}_2),\;
    \mat{\Sigma}_{11}-\mat{\Sigma}_{12}\mat{\Sigma}_{22}^{-1}\mat{\Sigma}_{21}\right)$.
  \end{itemize}
\end{itemize}
\end{minipage}}

% ***************************** question 1 **************************

\begin{question}{20}
  Each part needs only one or two of the results in the box.
  \begin{qparts}
    \qpart[2]{Evaluate $\displaystyle\int_{-\infty}^{\infty}\e{-x^2/2+3x}\,\dd x$.
      \begin{answer}
        Result 1 with $a=1$, $b=3$:
        $\sqrt{2\pi}\;\e{9/2}\approx 225.6$.
      \end{answer}}{\lines{3}}

    \qpart[3]{$X\sim\mathcal{N}(1,4)$.  Use the moment generating
      function to find $\av{\e{X}}$, and $\av{X^2}$.
      \begin{answer}
        $\av{\e{tX}} = \e{t+2t^2}$, so at $t=1$, $\av{\e{X}} = \e{3}$.
        Differentiating twice at $t=0$, or using
        $\av{X^2} = \mu^2+\sigma^2$, gives $\av{X^2} = 1+4 = 5$.
        (Note $\av{\e{X}}\neq\e{\av{X}}=\e{1}$.)
      \end{answer}}{\lines{4}}

    \qpart[3]{$X\sim\mathcal{N}(1,1)$ and $Y\sim\mathcal{N}(2,3)$ are
      independent.  What are the distributions of $X+Y$ and of $2X-Y$?
      \begin{answer}
        $X+Y\sim\mathcal{N}(3,4)$.  For $2X-Y$, the mean is
        $2\times1-2 = 0$ and the variance is $2^2\times1+(-1)^2\times3 = 7$,
        so $2X-Y\sim\mathcal{N}(0,7)$.  (Variances add even when we
        subtract.)
      \end{answer}}{\lines{4}}

    \qpart[3]{$\bm{x}\sim\mathcal{N}(\bm{0},\mat{I})$ in two dimensions
      and $\bm{z} = \mat{W}\bm{x}$ with
      $\mat{W} = \begin{pmatrix}1&1\\1&-1\end{pmatrix}$.  What is the
      distribution of $\bm{z}$?  Are $z_1 = x_1+x_2$ and $z_2 = x_1-x_2$
      independent?
      \begin{answer}
        $\bm{z}\sim\mathcal{N}(\bm{0},\mat{W}\mat{W}^\tr)$ with
        $\mat{W}\mat{W}^\tr = \begin{pmatrix}2&0\\0&2\end{pmatrix}$.  The
        covariance is diagonal, and for jointly Gaussian variables
        uncorrelated means independent, so yes.
      \end{answer}}{\lines{4}}

    \qpart[2]{$\bm{z}\sim\mathcal{N}(\bm{\mu},\mat{\Sigma})$ with
      $\bm{\mu} = (1,2,3)^\tr$ and
      $\mat{\Sigma} = \begin{pmatrix}4&1&0\\1&3&1\\0&1&2\end{pmatrix}$.
      What is the marginal distribution of $(z_1,z_3)$?
      \begin{answer}
        Strike out the second row and column:
        $(z_1,z_3)\sim\mathcal{N}\!\left((1,3)^\tr,
        \begin{pmatrix}4&0\\0&2\end{pmatrix}\right)$, so $z_1$ and $z_3$
        are independent, even though both are correlated with $z_2$.
      \end{answer}}{\lines{4}}

    \qpart[3]{$(z_1,z_2)$ is Gaussian with mean $\bm{0}$ and covariance
      $\begin{pmatrix}2&1\\1&1\end{pmatrix}$.  We observe $z_2 = 2$.
      What is the distribution of $z_1$ now?  Compare it with the
      distribution of $z_1$ before the observation.
      \begin{answer}
        Mean $0 + 1\times1^{-1}\times(2-0) = 2$; variance
        $2 - 1\times1^{-1}\times1 = 1$.  So $z_1\mid z_2=2\sim\mathcal{N}(2,1)$:
        the observation moved the mean from 0 to 2 and halved the
        variance (from 2).
      \end{answer}}{\lines{5}}

    \qpart[4]{Our prior belief about a quantity is $x\sim\mathcal{N}(0,4)$.
      We measure $y = x+\epsilon$ with noise $\epsilon\sim\mathcal{N}(0,1)$
      and get $y=3$.  Write down the joint distribution of $(x,y)$ and
      hence find the distribution of $x$ given $y=3$.
      \begin{answer}
        $y$ has mean 0 and variance $4+1=5$, and
        $\mathrm{Cov}(x,y) = \mathrm{Var}(x) = 4$, so
        $(x,y)\sim\mathcal{N}\!\left(\bm{0},
        \begin{pmatrix}4&4\\4&5\end{pmatrix}\right)$.  Conditioning,
        $\av{x\mid y} = \frac45\times3 = 2.4$ and
        $\mathrm{Var}(x\mid y) = 4-\frac{16}{5} = 0.8$.  The estimate is
        pulled most of the way to the measurement, and
        $1/0.8 = 1/4+1/1$: precisions add.
      \end{answer}}{\lines{6}}
  \end{qparts}
\end{question}
\clearpage

% ***************************** question 2 **************************

\begin{question}{15}
  A Gaussian process with covariance function
  $K(x,x') = \e{-(x-x')^2/2}$ has, for observations $\bm{y}$ at inputs
  $x_1,\ldots,x_m$ with noise variance $\sigma^2$, the predictive
  distribution at a new input $x_*$
  \begin{align*}
    f_*\mid\bm{y}\sim\mathcal{N}\!\left(\bm{k}_*^\tr(\mat{K}+\sigma^2\mat{I})^{-1}\bm{y},\;
    k_* - \bm{k}_*^\tr(\mat{K}+\sigma^2\mat{I})^{-1}\bm{k}_*\right)
  \end{align*}
  where $K_{ij} = K(x_i,x_j)$, $(\bm{k}_*)_i = K(x_i,x_*)$ and
  $k_* = K(x_*,x_*)$.
  \begin{qparts}
    \qpart[3]{Explain why this is the conditional result in the box,
      with $\bm{z}_1 = f_*$ and $\bm{z}_2 = \bm{y}$.  What are
      $\mat{\Sigma}_{11}$, $\mat{\Sigma}_{12}$ and $\mat{\Sigma}_{22}$?
      \begin{answer}
        $f_*$ and the observations are jointly Gaussian with mean zero.
        $\mat{\Sigma}_{11} = k_*$ (the prior variance of $f_*$),
        $\mat{\Sigma}_{12} = \bm{k}_*^\tr$ (the covariance of $f_*$ with each
        observation) and $\mat{\Sigma}_{22} = \mat{K}+\sigma^2\mat{I}$ (the
        observations' covariance, with the noise variance added on the
        diagonal).  Substituting into the conditional formula gives the
        result.
      \end{answer}}{\lines{5}}

    \qpart[4]{Take a single noise-free observation ($\sigma^2=0$),
      $y_1 = 1$ at $x_1 = 0$.  Find the predictive mean and variance at
      $x_*$.  What are they at $x_*=0$, and as $x_*\to\infty$?
      \begin{answer}
        $\mat{K} = 1$ and $k_* = 1$, and $\bm{k}_* = \e{-x_*^2/2}$.  So the mean
        is $\e{-x_*^2/2}$ and the variance is $1-\e{-x_*^2}$.  At $x_*=0$
        the mean is 1 and the variance 0: the prediction passes through
        the data point exactly.  As $x_*\to\infty$ the mean tends to 0
        and the variance to 1, the prior: far from the data we know
        nothing more than before.
      \end{answer}}{\lines{6}}

    \qpart[3]{Repeat with noise $\sigma^2 = 1$.  What changes at $x_* = 0$?
      \begin{answer}
        Now the mean is $\e{-x_*^2/2}/2$ and the variance
        $1-\e{-x_*^2}/2$.  At $x_*=0$ the mean is $\frac12$ and the
        variance $\frac12$: a noisy observation is only partly
        trusted (this is the noisy-measurement result of Question~1(g)).
      \end{answer}}{\lines{5}}

    \qpart[5]{Now take two noise-free observations, $y=1$ at $x=-1$ and
      $y=1$ at $x=1$.  Find the predictive mean and variance at $x_*=0$.
      (Hint: $\bm{y}=(1,1)^\tr$ is an eigenvector of $\mat{K}$.)
      \begin{answer}
        $\mat{K} = \begin{pmatrix}1&\e{-2}\\ \e{-2}&1\end{pmatrix}$ and
        $\bm{k}_* = \e{-1/2}(1,1)^\tr$.  Since
        $\mat{K}(1,1)^\tr = (1+\e{-2})(1,1)^\tr$,
        $\mat{K}^{-1}(1,1)^\tr = (1,1)^\tr/(1+\e{-2})$.  So the mean is
        $2\e{-1/2}/(1+\e{-2}) \approx 1.068$ and the variance is
        $1 - 2\e{-1}/(1+\e{-2})\approx 0.352$.  Between two points the mean
        overshoots 1 slightly, and the uncertainty is smaller than with one
        point but not zero.
      \end{answer}}{\lines{6}}
  \end{qparts}
\end{question}
\clearpage

% ***************************** question 3 **************************

\begin{question}{20}
  This question derives the normalisation of the multivariate normal
  step by step.
  \begin{qparts}
    \qpart[5]{Consider the integral
      \begin{align*}
        I_1 = \int_{-\infty}^{\infty} \e{-x^2/2} \, \dd x.
      \end{align*}
      Directly evaluating this is difficult, but there is a trick.
      Consider instead
      \begin{align*}
        I_1^2 &=  \int_{-\infty}^{\infty} \e{-x^2/2} \, \dd x
                \int_{-\infty}^{\infty} \e{-y^2/2} \, \dd y
                =   \int_{-\infty}^{\infty}  \int_{-\infty}^{\infty}
                  \e{-(x^2+y^2)/2} \, \dd x \, \dd y.
      \end{align*}
      By making the change of variables to polar coordinates where $r
      = \sqrt{x^2 + y^2}$ and $\theta = \arctan(y/x)$  (so that
      $x=r\,\cos(\theta)$, $y=r\,\sin(\theta)$) then  $\dd
      x\,\dd y = r\,\dd r\, \dd \theta$. Note that to integrate over
      all space we let $\theta$ vary from 0 to $2\pi$ and $r$ to vary
      from 0 to $\infty$.  Write down the integral in
      polar coordinates, make the further change of variables
      $u=r^2/2$ to evaluate $I_1^2$, and hence compute $I_1$.
    
\begin{answer}
      \begin{align*}
        I_1^2 = \int_0^{2\,\pi} \dd \theta \int_{0}^\infty r\,e^{-r^2/2} \, \dd r
        = 2\,\pi \ \int_{0}^\infty r\,e^{-r^2/2} \, \dd r.
      \end{align*}
      Making the change of variables $u=r^2/2$ so that $\dd u/\dd r =
      r$ or $\dd u= r\,\dd r$ we get 
      \begin{align*}
        I_1^2 = 2\,\pi\ \int_{0}^\infty e^{-u} \, \dd u = 2\,\pi
      \end{align*}
      so $I_1 = \sqrt{2\,\pi}$.
    \end{answer}
}{\lines{6}}
    \qpart[5]{By making a change of variables compute
      \begin{align*}
        I_2 = \int_{-\infty}^{\infty} \e{-(x-\mu)^2/(2\sigma^2)} \dd x
      \end{align*}
\begin{answer}
      The trick is to let $z=(x-\mu)/\sigma$ so that $\dd z = \dd
      x/\sigma$ or $\dd x = \sigma \dd z$.  Then
      \begin{align*}
        I_2 = \sigma \int_{-\infty}^{\infty} \e{-z^2/2} \dd z =
        \sigma\,I_1 = \sqrt{2\,\pi}\,\sigma
      \end{align*}
      Note that the \textit{probability density function} (PDF) for a
      normally distributed random variable is given by
      \begin{align*}
        \mathcal{N}(x\mid \mu, \sigma^2) = \frac{1}{\sqrt{2\,\pi}\,\sigma} \e{-(x-\mu)^2/(2\sigma^2)} .
      \end{align*}
      (Observe that at its peak $\mathcal{N}(\mu\mid \mu, \sigma^2) =
      1/(\sqrt{2\,\pi}\,\sigma)$, so when $\sigma<1/\sqrt{2\,\pi}$ the
      density exceeds 1, showing that PDFs are not probabilities.)
    \end{answer}
}{\lines{4}}
    \newpage
    
    \qpart[3]{By using the identity $\e{a+b} = \e{a}\,\e{b}$, or more
      generally
      \begin{align*}
        \e{\sum_i a_i} = \prod_i \e{a_i}
      \end{align*}
      compute
      \begin{align*}
        I_3 = \int_{-\infty}^\infty \cdots \int_{-\infty}^\infty
        \e{-\frac{1}{2}\|\bm{x}\|_2^2} \dd x_1 \cdots \dd x_n
      \end{align*}
      where $\bm{x} = (x_1, x_2, \ldots, x_n)^\tr$.
\begin{answer}
        We note that $\|\bm{x}\|_2^2 = \sum\limits_{i=1}^n x_i^2$
        \begin{align*}
          I_3 &= \int_{-\infty}^\infty \cdots \int_{-\infty}^\infty
        \e{-\frac{1}{2}  \sum\limits_{i=1}^n x_i^2} \dd x_1 \cdots \dd
                x_n \\
          &= \prod_{i=1}^n \int_{-\infty}^\infty  \e{-x_i^2/2} \,
            \dd x_i =  \prod_{i=1}^n \sqrt{2\,\pi} = (2\,\pi)^{n/2}.
        \end{align*}
      \end{answer}
}{\lines{6}}

      \qpart[6]{By using the fact that for a positive definite
        matrix, $\mat{\Sigma}$, we can use the eigenvector decomposition
        $\mat{\Sigma}^{-1} = \mat{V}\,\mat{\Lambda}^{-1}\,\mat{V}^\tr$ where
        $\mat{V}$ is an orthogonal matrix with determinant
        $\det{\mat{V}}=\pm1$ and $\mat{\Lambda}^{-1}$ is a diagonal matrix with
        elements $\lambda_i^{-1}$ compute
      \begin{align*}
        I_4 =  \int_{-\infty}^\infty \cdots \int_{-\infty}^\infty
        \e{-\frac{1}{2}  (\bm{x}-\bm{\mu})^\tr \mat{\Sigma}^{-1} (\bm{x}-\bm{\mu})} \, \dd x_1 \cdots \dd x_n.
      \end{align*}
\begin{answer}
        This needs some confidence to push through to the end.
        \begin{align*}
          I_4 &=  \int_{-\infty}^\infty \cdots \int_{-\infty}^\infty
          \e{-\frac{1}{2}  (\bm{x}-\bm{\mu})^\tr 
                 \mat{V}\,\mat{\Lambda}^{-1}\,\mat{V}^\tr (\bm{x}-\bm{\mu})}
                 \, \dd x_1 \cdots \dd x_n.
        \end{align*}
        We make the change of variables $\bm{y}= \mat{V}^\tr(\bm{x}-\bm{\mu})$.
        However, to make a change of variables,
        $\bm{x}\rightarrow\bm{y}$, where we use
        \begin{align*}
           \iiint_{\mathcal{R}(\bm{x})}
          f(\bm{x}) \, \dd x_1 \cdots \dd x_n = \iiint_{\mathcal{R}(\bm{y})}
          f(\bm{x}(\bm{y})) \, |\mat{J}|\, \dd y_1 \cdots \dd y_n 
        \end{align*}
        where $\mathcal{R}(\bm{y})$ is the same region as
        $\mathcal{R}(\bm{x})$ but specified in $\bm{y}$-coordinates
        and $\mat{J}$ is the Jacobian matrix with elements
        $J_{ij} = \partial x_i/\partial y_j$.  In our case
        $\bm{x}(\bm{y}) = \mat{V}\,\bm{y} + \bm{\mu}$
        \begin{align*}
          J_{ij} = \frac{\partial x_i}{\partial y_j} = \frac{\partial
          \ }{\partial y_j} \left(\sum_{k=1}^n V_{ik}y_k + \mu_i\right) = V_{ij}
        \end{align*}
        so that $\mat{J} = \mat{V}$ and $|\det{\mat{J}}| =
        |\det{\mat{V}}| = 1$.  This makes sense as the matrix
        $\mat{V}$ corresponds to a rotation (with a possible
        reflection) which does not change the volume.

        Thus on making the change of variables
        \begin{align*}
           I_4 &=  \int_{-\infty}^\infty \cdots \int_{-\infty}^\infty
          \e{-\frac{1}{2}  \bm{y}^\tr \mat{\Lambda}^{-1} \bm{y}}
                  \, \dd y_1 \cdots \dd y_n.
                  = \prod_{i=1}^n  \int_{-\infty}^\infty
                  \e{-y_i^2/(2\,\lambda_i)} \, \dd y_i = \prod_i
                 \sqrt{2\,\pi\,\lambda_i}
        \end{align*}
        where we used $I_2$ with $\sigma = \sqrt{\lambda_i}$.  (If
        $\mat{\Sigma}$ were not positive definite some $\lambda_i\leq0$
        and the integral would not converge.)
      \end{answer}
}{\lines{9}}
      \freshpage
      
      \qpart[1]{Using the facts, that $\mat{\Sigma} =
        \mat{V}\,\mat{\Lambda}\,\mat{V}^\tr$, for any two square
        matrices $\mat{A}$ and $\mat{B}$ the determinants satisfy
        $\det{\mat{A}\,\mat{B}} = \det{\mat{A}}\,\det{\mat{B}}$, and
        $\det{\mat{V}}=\det{\mat{V}^\tr}=\pm{1}$ show that $\det{\mat{\Sigma}}=\prod_i
        \lambda_i$.
\begin{answer}
        \begin{align*}
          \det{\mat{\Sigma}} = \det{\mat{V}\,\mat{\Lambda}\,\mat{V}^\tr} =
          \det{\mat{V}}\,\det{\mat{\Lambda}}\,\det{\mat{V}^\tr}
          = \det{\mat{\Lambda}} = \prod_{i=1}^n \lambda_i
        \end{align*}
        We have used that $\det{\mat{V}}=\det{\mat{V}^\tr}$ which equals 1 or
        -1, but $\det{\mat{V}}\times\det{\mat{V}^\tr}=1$.  Also we use that
        the determinant of a diagonal matrix is equal to the product
        of the diagonal elements.  We note that
        \begin{align*}
          I_4 = (2\,\pi)^{n/2} \sqrt{\det{\mat{\Sigma}}}
        \end{align*}
        There is another trick to simplify the notation a bit more.
        For an $n\times n$ matrix $\mat{M}$ then
        $\det{c\,\mat{M}}=c^n\,\det{\mat{M}}$ as the determinant of the matrix
        involves a sum of terms where each term involves a product of
        $n$ elements.  Multiplying each element by $c$ means that the
        determinant increases by $c^n$.  Thus, we can write
        \begin{align*}
            I_4 = (2\,\pi)^{n/2} \sqrt{\det{\mat{\Sigma}}}= \sqrt{\det{2\,\pi\,\mat{\Sigma}}}
        \end{align*}
        and the multivariate normal PDF is
        \begin{align*}
          \mathcal{N}(\bm{x} \mid \bm{\mu}, \mat{\Sigma}) =
          \frac{1}{\sqrt{\det{2\,\pi\,\mat{\Sigma}}}}  \, \e{-\frac{1}{2}  (\bm{x}-\bm{\mu})^\tr \mat{\Sigma}^{-1} (\bm{x}-\bm{\mu})}
        \end{align*}
      \end{answer}
}{\lines{6}}
    \end{qparts}
  \end{question}
  \clearpage
% ***************************** question 4 **************************


\begin{question}{25}
  This question is harder: it derives the marginal and conditional
  results in the box by doing the integrals directly (the lecture
  used a shorter route).  Take it slowly; each part is a short
  calculation.

  Consider a multivariate normal distribution
  \begin{align*}
    f_{\bm{X},\bm{Y}}(\bm{x}, \bm{y}) = \mathcal{N}
    \!\left( \begin{pmatrix}\bm{x}\\ \bm{y}\end{pmatrix}
    \,\middle|\, \begin{pmatrix} \bm{0} \\
      \bm{0} \end{pmatrix}, \begin{pmatrix} \mat{A} & \mat{B}\\
      \mat{B}^{\tr} & \mat{C} \end{pmatrix} \right) 
  \end{align*}
  where $\mat{A}$ and $\mat{C}$ are symmetric (positive definite)
  matrices.  The matrix
  \begin{align*}
    \mat{\Sigma} &= \begin{pmatrix} \mat{A} & \mat{B}\\ \mat{B}^{\tr} & \mat{C} \end{pmatrix}
  \end{align*}
  is the covariance matrix.

  We want to compute the conditional probability density
  function $f_{\bm{X},\bm{Y}}(\bm{x}\mid \bm{y})$.  This is complicated because the
  normal distribution involves the inverse of the covariance matrix.  Let
      \begin{align*}
        \mat{U} &= \begin{pmatrix} \Id & \mat{B}\,\mat{C}^{-1}\\
          \mat{0} & \Id \end{pmatrix},
         & \mat{D} &=  \begin{pmatrix} \mat{A}-
           \mat{B}\,\mat{C}^{-1}\mat{B}^{\tr}
           & \mat{0}\\ \mat{0}^{\tr} & \mat{C}  \end{pmatrix}
      \end{align*}
      where $\Id$ is the identity matrix.
  \begin{qparts}
    \qpart[3]{Compute $\mat{U}\,\mat{D}$
\begin{answer}
      \begin{align*}
        \mat{U}\,\mat{D}&=\begin{pmatrix} \Id &
          \mat{B}\,\mat{C}^{-1}\\ \mat{0} &
          \Id \end{pmatrix} \begin{pmatrix} \mat{A}-
          \mat{B}\,\mat{C}^{-1}\mat{B}^{\tr} & \mat{0}\\ \mat{0}^{\tr}
          & \mat{C} \end{pmatrix}\\
        &=\begin{pmatrix} \mat{A}- \mat{B}\,\mat{C}^{-1}\mat{B}^{\tr}
          & \mat{B}\\ \mat{0}^{\tr} & \mat{C} \end{pmatrix}
      \end{align*}
    \end{answer}
}{\lines{5}}
    
    \qpart[3]{Using the previous result compute
      $(\mat{U}\,\mat{D})\,\mat{U}^\tr$.  Hence show
      $\mat{\Sigma} = \mat{U}\,\mat{D}\,\mat{U}^\tr$.
\begin{answer}
      \begin{align*}
        (\mat{U}\,\mat{D})\,\mat{U}^{\tr}&=\begin{pmatrix} \mat{A}-
          \mat{B}\,\mat{C}^{-1}\mat{B}^{\tr} & \mat{B}\\ \mat{0}^{\tr} &
          \mat{C} \end{pmatrix}\begin{pmatrix} \Id & \mat{0}\\
          \mat{C}^{-1}\mat{B}^{\tr} & \Id \end{pmatrix} =\begin{pmatrix}
          \mat{A} & \mat{B}\\ \mat{B}^{\tr} & \mat{C} \end{pmatrix} = \mat{\Sigma}
      \end{align*}
    \end{answer}
}{\lines{5}}
    \clearpage
    
    \qpart[1]{Given that $\mat{\Sigma} = \mat{U}\,\mat{D}\,\mat{U}^\tr$
      write down $\mat{\Sigma}^{-1}$ in terms of $\mat{U}$ and
      $\mat{D}$
\begin{answer}
      \begin{align*}
        \mat{\Sigma}^{-1} = (\mat{U}^{\tr})^{-1}\mat{D}^{-1} \mat{U}^{-1}
      \end{align*}
    \end{answer}
}{\lines{1}}
    
    \qpart[6]{Demonstrate by direct multiplication that
      \begin{align*}
        \mat{U}^{-1}
        &= \begin{pmatrix} \Id &
          -\mat{B}\,\mat{C}^{-1} \\ \mat{0} & \Id \end{pmatrix}
                                   &
         \mat{D}^{-1}
         &=\begin{pmatrix} (\mat{A}-
           \mat{B}\,\mat{C}^{-1}\mat{B}^{\tr})^{-1} & \mat{0}\\
           \mat{0}^{\tr} & \mat{C}^{-1} \end{pmatrix} 
      &
      (\mat{U}^\tr)^{-1} = \begin{pmatrix} \Id &
        \mat{0} \\ -\mat{C}^{-1} \mat{B}^\tr & \Id \end{pmatrix} 
      \end{align*}
      i.e. show $\mat{U}^{-1}\mat{U}=\Id$,
      $\mat{D}^{-1}\mat{D}=\Id$ and
      $(\mat{U}^\tr)^{-1}\mat{U}^\tr= \Id$.
    
\begin{answer}
    \begin{enumerate}
      \item
        \begin{align*}
          \mat{U}^{-1}\mat{U} = \begin{pmatrix} \Id &
            -\mat{B}\,\mat{C}^{-1} \\ \mat{0} &
            \Id \end{pmatrix} \begin{pmatrix} \Id &
            \mat{B}\,\mat{C}^{-1}\\ \mat{0} & \Id \end{pmatrix}
          = \begin{pmatrix} \Id & \mat{0} \\ \mat{0} &
            \Id \end{pmatrix} 
        \end{align*}
      \item
        \begin{align*}
          \mat{D}^{-1}\mat{D} =  \begin{pmatrix} (\mat{A}-
           \mat{B}\,\mat{C}^{-1}\mat{B}^{\tr})^{-1} & \mat{0}\\
           \mat{0}^{\tr} & \mat{C}^{-1} \end{pmatrix} \begin{pmatrix} \mat{A}-
           \mat{B}\,\mat{C}^{-1}\mat{B}^{\tr}
           & \mat{0}\\ \mat{0}^{\tr} &
           \mat{C}  \end{pmatrix}= \begin{pmatrix} \Id & \mat{0}
           \\ \mat{0} &  \Id \end{pmatrix}                                                                                   
        \end{align*}
        where we use $(\mat{A}-
           \mat{B}\,\mat{C}^{-1}\mat{B}^{\tr})^{-1}(\mat{A}-
           \mat{B}\,\mat{C}^{-1}\mat{B}^{\tr}) = \Id$ and
           $\mat{C}^{-1}\mat{C}=\Id$.
         \item
           \begin{align*}
              (\mat{U}^\tr)^{-1}\mat{U}^\tr =
           \begin{pmatrix} \Id &
        \mat{0} \\ -\mat{C}^{-1} \mat{B}^\tr & \Id \end{pmatrix} 
         \begin{pmatrix} \Id & \mat{0}\\ 
           \mat{C}^{-1}\mat{B}^{\tr} & \Id \end{pmatrix}
                                       =          \begin{pmatrix} \Id & \mat{0}
           \\ \mat{0} &  \Id \end{pmatrix}                                             
           \end{align*}
         \end{enumerate}
       \end{answer}
}{\inumberedlines[3]{3}}
       \clearpage
       
       \qpart[4]{Letting $\bm{z} = \binom{\bm{x}}{\bm{y}}$ then we can
         write
         \begin{align*}
           f_{\bm{X},\bm{Y}}(\bm{x}, \bm{y}) = f_{\bm{Z}}(\bm{z}) = \mathcal{N}(\bm{z}\mid\bm{0},
           \mat{\Sigma}) = \frac{1}{\sqrt{\det{2\,\pi\,\mat{\Sigma}}}} \,
           \e{- \frac{1}{2}\, \bm{z}^{\tr}\mat{\Sigma}^{-1} \bm{z}}.
         \end{align*}
         where $\det{2\,\pi\,\mat{\Sigma}}$ is the determinant of the matrix
         $2\,\pi\,\mat{\Sigma}$ and is introduced to ensure that
         $f_{\bm{X},\bm{Y}}(\bm{x},\bm{y})$ is normalised.  From parts (c) and (d)
         \begin{align*}
           \mat{\Sigma}^{-1} = \begin{pmatrix} \Id & \mat{0} \\
             -\mat{C}^{-1} \mat{B}^\tr&
             \Id \end{pmatrix} \begin{pmatrix} (\mat{A}-
             \mat{B}\,\mat{C}^{-1}\mat{B}^{\tr})^{-1} & \mat{0}\\
             \mat{0}^{\tr} &
             \mat{C}^{-1} \end{pmatrix} \begin{pmatrix} \Id &
             -\mat{B}\,\mat{C}^{-1} \\ \mat{0} &
             \Id \end{pmatrix}. 
         \end{align*}
         Expand out $\bm{z}^{\tr}\mat{\Sigma}^{-1} \bm{z} =
         (\bm{x}^\tr, \bm{y}^\tr)\mat{\Sigma}^{-1}
         \binom{\bm{x}}{\bm{y}}$ (start by multiplying the vectors
         $\bm{z}^\tr$ by $(\mat{U}^\tr)^{-1}$ and $\bm{z}$ by $\mat{U}^{-1}$)
       
\begin{answer}
         \begin{align*}
           \bm{z}^{\tr}\mat{\Sigma}^{-1}\bm{z} &=
                                              (\bm{x}^{\tr},\bm{y}^{\tr})\begin{pmatrix}
                                                \Id & \mat{0} \\
                                                -\mat{C}^{-1}
                                                \mat{B}^{\tr} &
                                                \Id \end{pmatrix} \begin{pmatrix}
                                                (\mat{A}-
                                                \mat{B}\,\mat{C}^{-1}\mat{B}^\tr)^{-1}
                                                & \mat{0}\\
                                                \mat{0}^{\tr} &
                                                \mat{C}^{-1} \end{pmatrix} \begin{pmatrix}
                                                \Id &
                                                -\mat{B}\,\mat{C}^{-1}
                                                \\ \mat{0} &
                                                \Id \end{pmatrix} \begin{pmatrix}\bm{x}\\
                                                \bm{y}\end{pmatrix}\\ 
                                            &=
                                              (\bm{x}^{\tr}-\bm{y}^\tr\mat{C}^{-1}\mat{B}^{\tr},\,\bm{y}^{\tr}) \begin{pmatrix}
                                                (\mat{A}-
                                                \mat{B}\,\mat{C}^{-1}\mat{B}^{\tr})^{-1}
                                                & \mat{0}\\
                                                \mat{0}^{\tr} &
                                                \mat{C}^{-1} \end{pmatrix} \begin{pmatrix}\bm{x}-\mat{B}\,\mat{C}^{-1}\bm{y}\\
                                                \bm{y}\end{pmatrix}\\ 
                                            &=
                                              (\bm{x}^{\tr}-\bm{y}^{\tr}\mat{C}^{-1}\mat{B}^{\tr})(\mat{A}-
                                              \mat{B}\,\mat{C}^{-1}\mat{B}^{\tr})^{-1}(\bm{x}-\mat{B}\,\mat{C}^{-1}\bm{y})
                                              + \bm{y}^{\tr} \mat{C}^{-1}
                                              \bm{y}
         \end{align*}
       \end{answer}
}{\lines{10}}
       \clearpage
       
       \qpart[5]{(We write $\int f(\bm{z})\,\dd\bm{z}$ for
         $\int\cdots\int f(\bm{z})\,\dd z_1\cdots\dd z_n$.)  To compute $f_{\bm{X}\mid\bm{Y}}(\bm{x}\mid\bm{y}) = f_{\bm{X},\bm{Y}}(\bm{x},
         \bm{y})/f_{\bm{Y}}(\bm{y})$ we need to find
         \begin{align*}
           f_{\bm{Y}}(\bm{y})&= \int_{-\infty}^{\infty} f(\bm{x}, \bm{y}) \,\dd \bm{x}\\
                      &= \int_{-\infty}^{\infty}
                        \frac{1}{\sqrt{\det{2\,\pi\,\mat{\Sigma}}}} \,
                        \e{- \frac{1}{2}\,
                        (\bm{x}^{\tr}-\bm{y}^{\tr}\mat{C}^{-1}\mat{B}^{\tr})(\mat{A}-
                        \mat{B}\,\mat{C}^{-1}\mat{B}^\tr)^{-1}(\bm{x}
                        -\mat{B}\,\mat{C}^{-1}\bm{y}) - \frac{1}{2}\,
                        \bm{y}^{\tr} \mat{C}^{-1} \bm{y}} \dd \bm{x} \\
                      &=  \frac{\e{- \frac{1}{2}\,  \bm{y}^{\tr} \mat{C}^{-1} \bm{y}}}{\sqrt{\det{2\,\pi\,\mat{\Sigma}}}}
                        \int_{-\infty}^{\infty}
                        \e{- \frac{1}{2}\,
                        (\bm{x}^{\tr}-\bm{y}^{\tr}\mat{C}^{-1}\mat{B}^{\tr})(\mat{A}-
                        \mat{B}\,\mat{C}^{-1}\mat{B}^\tr)^{-1}(\bm{x}
                        -\mat{B}\,\mat{C}^{-1}\bm{y}) } \dd \bm{x} 
         \end{align*}
         By making a change of variable from $\bm{x}$ to $\bm{u} =
         \bm{x}-\mat{B}\,\mat{C}^{-1}\bm{y}$
         rewrite the integral
         \begin{align*}
           I = \int_{-\infty}^{\infty}
           \e{- \frac{1}{2}\,
           (\bm{x}^{\tr}-\bm{y}^{\tr}\mat{C}^{-1}\mat{B}^{\tr})(\mat{A}-
           \mat{B}\,\mat{C}^{-1}\mat{B}^\tr)^{-1}(\bm{x}
           -\mat{B}\,\mat{C}^{-1}\bm{y}) } \dd \bm{x} 
         \end{align*}
         then use
         \begin{align*}
           \int_{-\infty}^{\infty}
           \e{- \frac{1}{2}\, \bm{u}^\tr \mat{M}^{-1} \,
           \bm{u}} \dd \bm{u} = \sqrt{\det{2\,\pi\,\mat{M}}}  
         \end{align*}
         to evaluate $f_{\bm{Y}}(\bm{y})$.
       
\begin{answer}
         Making the change of variables $\bm{u} =
         \bm{x}-\mat{B}\,\mat{C}^{-1}\bm{y}$
          \begin{align*}
           I  =   \int_{-\infty}^{\infty}
           \e{- \frac{1}{2}\,\bm{u}^\tr (\mat{A}-
           \mat{B}\,\mat{C}^{-1}\mat{B}^\tr)^{-1}\bm{u}} \, \dd \bm{u}
            = \sqrt{\det{2\,\pi\,(\mat{A}-\mat{B}\,\mat{C}^{-1}\mat{B}^\tr)}}  
         \end{align*}
         Thus
         \begin{align*}
            f_{\bm{Y}}(\bm{y})&=\sqrt{\frac{\det{2\,\pi\,(\mat{A}-
                         \mat{B}\,\mat{C}^{-1}\mat{B}^\tr)}}{\det{2\,\pi\,\mat{\Sigma}}}}
                         \, \e{- \frac{1}{2}\, \bm{y}^{\tr} \mat{C}^{-1}
                         \bm{y}}.
         \end{align*}
         where $\det{\mat{\Sigma}}$ denotes the determinant of a matrix $\mat{\Sigma}$.
         In passing we recall that $\mat{\Sigma}=\mat{U}\mat{D}\mat{U}^\tr$
         so $\det{\mat{\Sigma}}=\det{\mat{U}}\,\det{\mat{D}}\,\det{\mat{U}^\tr}$,  But it
         is easy to show $\det{\mat{U}^{\tr}} = \det{\mat{U}} =1$  and because
         $\mat{D}$ is block diagonal $\det{\mat{D}} = \det{\mat{A}-
         \mat{B}\,\mat{C}^{-1}\mat{B}^{\tr}} \, \det{\mat{C}}$ so that  $\det{\mat{\Sigma}} = \det{\mat{A}-
         \mat{B}\,\mat{C}^{-1}\mat{B}^{\tr}} \, \det{\mat{C}}$.  Thus
         \begin{align*}
           f_{\bm{Y}}(\bm{y}) = \frac{1}{\sqrt{\det{2\,\pi\,\mat{C}}}}\, \e{-
           \frac{1}{2}\, \bm{y}^{\tr} \mat{C}^{-1} \bm{y}} = \mathcal{N}(\bm{y}
           \mid \bm{0}, \mat{C}).
         \end{align*}
         Recall that
           \begin{align*}
             f_{\bm{X},\bm{Y}}(\bm{x}, \bm{y}) = \mathcal{N}
             \!\left( \begin{pmatrix}\bm{x}\\ \bm{y}\end{pmatrix}
             \,\middle|\, \begin{pmatrix} \bm{0} \\
               \bm{0} \end{pmatrix}, \begin{pmatrix} \mat{A} & \mat{B}\\
               \mat{B}^{\tr} & \mat{C} \end{pmatrix} \right) 
           \end{align*}
           so we see $\bm{y}$ is normally distributed with covariance
           $\mat{C}$ as might be expected.
         \end{answer}
}{\lines{10}}
         \clearpage
         \qpart[3]{Using  $f_{\bm{X}\mid\bm{Y}}(\bm{x}\mid\bm{y}) = f_{\bm{X},\bm{Y}}(\bm{x},
         \bm{y})/f_{\bm{Y}}(\bm{y})$  write down
         $f_{\bm{X}|\bm{Y}}(\bm{x}\mid\bm{y})$.
\begin{answer}
         \begin{align*}
           f_{\bm{X}|\bm{Y}}(\bm{x}\mid\bm{y}) &= \frac{f_{\bm{X},\bm{Y}}(\bm{x},
           \bm{y})}{f_{\bm{Y}}(\bm{y})}
           =
           \frac{\det{2\,\pi\,\mat{C}}^{\frac{1}{2}}}{\det{2\,\pi\,\mat{\Sigma}}^{\frac{1}{2}}}
           \,\e{- \frac{1}{2}\,
           (\bm{x}^{\tr}-\bm{y}^{\tr}\mat{C}^{-1}\mat{B}^{\tr})\,(\mat{A}-
           \mat{B}\,\mat{C}^{-1}\mat{B}^\tr)^{-1}(\bm{x}-\mat{B}\,\mat{C}^{-1}\bm{y})}\\ 
           &= \frac{1}{\det{2\,\pi\,(\mat{A}- \mat{B}\,\mat{C}^{-1}\mat{B}^\tr)}^{\frac{1}{2}}}
             \,\e{- \frac{1}{2}\,
             (\bm{x}^{\tr}-\bm{y}^{\tr}\mat{C}^{-1}\mat{B}^{\tr})\,(\mat{A}-
             \mat{B}\,\mat{C}^{-1}\mat{B}^\tr)^{-1}(\bm{x}-\mat{B}\,\mat{C}^{-1}\bm{y})}\\  
           &= \mathcal{N}(\bm{x} \mid \mat{B}\,\mat{C}^{-1}\bm{y},\,\, \mat{A}-  \mat{B}\,\mat{C}^{-1}\mat{B}^\tr). 
         \end{align*}
         This is the conditioning result we use in deriving the
         Bayesian update for a Gaussian Process.
       \end{answer}
}{\lines{10}}
    \end{qparts}
\end{question}



\end{document}




